% Acrónimos del TFM — se usan con \gls{clave}, \glspl{clave}, \Gls{clave}

% --- Movilidad y política ---
\newacronym[longplural={zonas de bajas emisiones},plural={ZBE}]{zbe}{ZBE}{zona de bajas emisiones}
\newacronym{nap}{NAP}{Punto de Acceso Nacional de Transporte Multimodal}
\newacronym{mitrams}{MITRAMS}{Ministerio de Transportes y Movilidad Sostenible}
\newacronym{emtval}{EMT Valencia}{Empresa Municipal de Transportes de Valencia}
\newacronym{emtmad}{EMT Madrid}{Empresa Municipal de Transportes de Madrid}
% --- Datos geográficos y transporte ---
\newacronym{osm}{OSM}{OpenStreetMap}
\newacronym{gtfs}{GTFS}{General Transit Feed Specification}
\newacronym{fossgis}{FOSSGIS}{Free and Open Source Software for Geographical Information Systems}
\newacronym{ltds}{LTDS}{London Travel Demand Survey}


% --- Herramientas de enrutado ---
\newacronym{osrm}{OSRM}{Open Source Routing Machine}
\newacronym{otp}{OTP}{OpenTripPlanner}
\newacronym{mld}{MLD}{Multi-Level Dijkstra}
\newacronym{ch}{CH}{Contraction Hierarchies}
% --- Modelos de aprendizaje automático ---
\newacronym{rf}{RF}{Random Forest}
\newacronym{xgboost}{XGBoost}{eXtreme Gradient Boosting}

% --- Dataset y modelos de elección modal ---
\newacronym{lpmc}{LPMC}{London Passenger Mode Choice}
\newacronym[longplural={modelos de utilidad aleatoria},plural={RUM}]{rum}{RUM}{modelo de utilidad aleatoria}
\newacronym{mnl}{MNL}{modelo logit multinomial}
\newacronym{iia}{IIA}{independencia de alternativas irrelevantes}
\newacronym{gbdt}{GBDT}{Gradient Boosting Decision Trees}
\newacronym[longplural={redes neuronales profundas},plural={DNN}]{dnn}{DNN}{red neuronal profunda}
\newacronym{gmpca}{GMPCA}{Geometric Mean Probability of Correct Alternative}

% --- Metodología ---
\newacronym{mvp}{MVP}{producto mínimo viable}

% --- Infraestructura y sistema ---
\newacronym{lfs}{LFS}{Large File Storage}
\newacronym{cors}{CORS}{Cross-Origin Resource Sharing}

% --- Cartografía ---
\newacronym{pnoa}{PNOA}{Plan Nacional de Ortofotografía Aérea}
